\newpage
\begin{trapbox}
    \textbf{\textcolor{CTrap}{易错点一：只有垂直就套等长旋转公式}}
    \begin{center}
\begin{tikzpicture}[x=0.65cm,y=0.65cm,>=stealth,every node/.style={font=\small},line cap=round]

\draw[very thick,blue!75!black] (0,0)--(3,0) node[midway,below] {$3$};
\draw[very thick,teal!75!black] (0,0)--(0,3) node[midway,left] {$3$};
\draw[orange!85!black,very thick,dashed] (0,3)--(0,5) node[midway,left] {$2$};
\draw (.35,0)--(.35,.35)--(0,.35);
\node[below left] at (0,0) {$B(0,0)$};
\node[right] at (3,0) {$P(3,0)$};
\fill[teal!75!black] (0,3) circle (2pt) node[right] {$Q_{\rm rot}(0,3)$：旋转终点};
\fill[orange!85!black] (0,5) circle (2pt) node[right] {$Q(0,5)$：只有垂直};
\node[anchor=west] at (5,1.2) {方向相同，长度不同};

\end{tikzpicture}
\end{center}

    \textbf{错在哪里：} 图中 $BP=3$，$BQ=5$，虽然垂直，但比真正的旋转终点多出长度 $2$。
    \textbf{正确替代：} 只有垂直时，可写 $(s,t)=k(-v,u)$；还须用等长条件才能确定 $k=\pm1$。

    \medskip
    \textbf{\textcolor{CTrap}{易错点二：绕 $B$ 旋转，却直接改 $P$ 的坐标}}
    \begin{center}
\begin{tikzpicture}[x=0.65cm,y=0.65cm,>=stealth,every node/.style={font=\small},line cap=round]

\draw[->,gray] (-4,0)--(6,0) node[right] {$x$};
\draw[->,gray] (0,-1.9)--(0,5.9) node[above] {$y$};
\draw[blue!75!black,very thick] (2,-1)--(5,3);
\draw[teal!75!black,very thick] (2,-1)--(-2,2);
\draw[orange!85!black,dashed,thick] (0,0)--(5,3);
\draw[orange!85!black,dashed,thick] (0,0)--(-3,5);
\node[below left] at (0,0) {$O$}; \node[below] at (2,-1) {$B(2,-1)$};
\node[right] at (5,3) {$P(5,3)$};
\node[left] at (-2,2) {$Q(-2,2)$};
\node[above left] at (-3,5) {$W(-3,5)$};
\node[anchor=west,teal!75!black] at (8.5,4.6) {实线：绕 $B$};
\node[anchor=west,orange!85!black] at (8.5,3.4) {虚线：绕 $O$};
\node[anchor=west] at (8.5,1.3) {$W$ 是另一种旋转};
\node[anchor=west] at (8.5,.2) {不是本题的终点};

\end{tikzpicture}
\end{center}

    \textbf{错在哪里：} 把 $P(5,3)$ 写成 $(-3,5)$，实际是绕原点 $O$ 转了直角；虚线与实线的旋转中心不同。
    \textbf{正确替代：} 绕 $B(2,-1)$，先取 $P-B=(3,4)$，逆时针变为 $(-4,3)$，再加 $B$ 得 $Q=(-2,2)$；顺时针得 $(6,-4)$。

    \medskip
    \textbf{\textcolor{CTrap}{易错点三：只考虑一种转向}}
    \textbf{错在哪里：} 垂直且等长仍有两个位置，不能仅凭示意图偏向某一侧就漏掉一支。
    \textbf{正确替代：} 对照首页 $Q_1,Q_2$：保留顺、逆时针两种结果，再按题目的象限、直线或其他位置条件筛选。

    \newpage
    \textbf{\textcolor{CTrap}{易错点四：交换而不变号，或交换后两个都变号}}
    \begin{center}
\begin{tikzpicture}[x=0.65cm,y=0.65cm,>=stealth,every node/.style={font=\small},line cap=round]

\draw[->,gray] (-3,0)--(4,0); \draw[->,gray] (0,-4)--(0,4.3);
\draw[blue!75!black,very thick] (0,0)--(3,2);
\draw[teal!75!black,very thick] (0,0)--(-2,3);
\draw[orange!85!black,thick,dashed] (0,0)--(2,3);
\draw[orange!85!black,thick,dashed] (0,0)--(-2,-3);
\draw (.42,.28)--(.14,.70)--(-.28,.42);
\node[below right] at (0,0) {$B$}; \node[right] at (3,2) {$P(3,2)$};
\node[above left] at (-2,3) {$(-2,3)$}; \node[above right] at (2,3) {$(2,3)$};
\node[below left] at (-2,-3) {$(-2,-3)$};
\node[anchor=west,teal!75!black] at (4.8,3) {变一个号：数量积 $0$};
\node[anchor=west] at (4.8,1.4) {只交换：数量积 $12$};
\node[anchor=west] at (4.8,-.2) {交换后都变号：数量积 $-12$};
\node[anchor=west] at (4.8,-1.8) {等长仍在，垂直已丢失};

\end{tikzpicture}
\end{center}

    \textbf{错在哪里：} 把 $(3,2)$ 改为 $(2,3)$ 或 $(-2,-3)$，长度仍相等，但数量积为 $12$ 或 $-12$，已经不垂直。
    \textbf{正确替代：} 逆时针 $(-v,u)$，顺时针 $(v,-u)$；数量积都是 $0$。用方向理解符号，不只背“交换”。

    \medskip
    \textbf{\textcolor{CTrap}{易错点五：有竖直线还强行使用斜率乘积}}
    \textbf{错在哪里：} $k_{BP}k_{BQ}=-1$ 要求两条直线的斜率都存在，竖直线没有斜率。
    \textbf{正确替代：} 对照方法四右图，$P(0,3)$ 对应的旋转点为 $(-3,0)$ 或 $(3,0)$；用坐标旋转直接求解，或将水平、竖直情况单独处理。

    \medskip
    \textbf{\textcolor{CTrap}{易错点六：把一个坐标关系当成全部条件}}
    \textbf{错在哪里：} 例3中 $|y_1|=1$ 只完成了消元，还没有确定曲线上的点和对应的 $Q$。
    \textbf{正确替代：} 继续代回 $x_1^2+16y_1^2=25$，得 $x_1=\pm3$；需要具体点时，再按转向确定 $y_2$。
\end{trapbox}
